\documentclass[11pt]{article}
\usepackage[margin=1in]{geometry}
\usepackage{amsmath,amsthm,amssymb}
\usepackage{hyperref}

\newtheorem{theorem}{Theorem}[section]
\newtheorem{lemma}[theorem]{Lemma}
\newtheorem{remark}[theorem]{Remark}

\title{\textbf{A Rayleigh Lower Bound for Normalized Inverse-GCD\\
Matrices on Prime-Pair and Multi-Representation Vectors}}
\author{Jonathan Robert Simons\\
Savannah, Georgia, USA\\
\texttt{simonsmedicalinnovations@gmail.com}}
\date{Corrected preprint \quad|\quad August 2026}

\begin{document}
\maketitle

\begin{center}
\fbox{\parbox{0.92\linewidth}{\centering\small
\textbf{Disclaimer.} Restricted Rayleigh bound on Goldbach test vectors;
not a full-spectrum floor. No Navier--Stokes claim.}}
\end{center}
\vspace{0.75em}

\begin{abstract}
Let $\widetilde Q_N(i,j)=1/(\gcd(i,j)\sqrt{ij})$ for $1\le i,j\le N$.
For distinct primes $p,q$ set $v=e_p-e_q$. We prove
\[
R(v)=\frac{v^\top\widetilde Q_N v}{\|v\|_2^2}
=\frac12\Big(\frac1{p^2}+\frac1{q^2}\Big)-\frac1{\sqrt{pq}}
\;>\;-\frac12.
\]
The same floor holds for multi-representation vectors
$v_k=\sum_{p\in\mathcal{R}(k)}(e_p-e_{k-p})$ by a nonnegative cross-term
identity. We do not claim a full-spectrum floor
$\lambda_{\min}(\widetilde Q_N)>-1/2$ (false for moderate $N$).
\end{abstract}

\section{Operator}

\[
\widetilde Q_N(p,p)=\frac1{p^2},\qquad
\widetilde Q_N(p,q)=\frac1{\sqrt{pq}}\quad(p\neq q\text{ primes}).
\]

\section{Single pair}

\begin{theorem}
For distinct primes $p,q$,
\[
R(e_p-e_q)=\frac12\Big(\frac1{p^2}+\frac1{q^2}\Big)-\frac1{\sqrt{pq}}\;>\;-\frac12.
\]
\end{theorem}

\begin{proof}
\[
R+\frac12
=\frac12+\frac1{2p^2}+\frac1{2q^2}-\frac1{\sqrt{pq}}
>\frac12-\frac1{\sqrt{pq}}.
\]
Since $pq\ge 2\cdot 3=6$, one has $1/\sqrt{pq}\le 1/\sqrt6<1/2$, hence $R+1/2>0$.
\end{proof}

\begin{remark}
The Rayleigh quotient may be negative (e.g.\ $(p,q)=(2,3)$ gives $R\approx -0.228$)
but remains strictly above $-1/2$. Full-spectrum $\lambda_{\min}(\widetilde Q_N)$
falls below $-1/2$ already for small $N$; the present bound is restricted to
Goldbach test vectors.
\end{remark}

\section{Multi-representation}

\begin{lemma}[Cross-term factorization]
For distinct primes $p\neq q$ and $r\neq s$,
\[
\widetilde Q(p,r)-\widetilde Q(p,s)-\widetilde Q(q,r)+\widetilde Q(q,s)
=\bigl(p^{-1/2}-q^{-1/2}\bigr)\bigl(r^{-1/2}-s^{-1/2}\bigr).
\]
If $p<q$ and $r<s$, the cross term is strictly positive.
\end{lemma}

\begin{theorem}[Multi-representation Bridge$^*$]
Let $k$ be even and
$\mathcal{R}(k)=\{p\text{ prime}:2\le p<k/2,\ k-p\text{ prime}\}$.
Set $v_k=\sum_{p\in\mathcal{R}(k)}(e_p-e_{k-p})\neq 0$. Then $R(v_k)>-1/2$.
\end{theorem}

\begin{proof}
Unordered Goldbach pairs for fixed $k$ are pairwise disjoint.
Write $v_k=\sum_a v_a$ with $v_a=e_{p_a}-e_{q_a}$, $p_a<q_a$.
Diagonal blocks give $R(v_a)>-1/2$ by the single-pair theorem; cross blocks
are positive by the lemma. Hence $R(v_k)\ge\min_a R(v_a)>-1/2$.
\end{proof}

Full operator note: \texttt{04\_q6\_inverse\_gcd.tex}.

\end{document}
